\documentclass[11pt,spanish]{article}

% =========================================================
% PLANTILLA BASE PARA INFORMES ACADÉMICOS / TÉCNICOS
% =========================================================

% --- Idioma y codificación ---
\usepackage[spanish,es-tabla]{babel}
\usepackage[T1]{fontenc}
\usepackage[utf8]{inputenc}
\usepackage{textcomp}
\DeclareUnicodeCharacter{20A1}{\textcolonmonetary{}}

% --- Márgenes / papel ---
\usepackage{geometry}
\geometry{a4paper, margin=2.7cm}

% --- Matemáticas ---
\usepackage{amsmath}

% --- Tablas, enlaces y elementos visuales ---
\usepackage{array}
\usepackage{tabularx}
\usepackage{booktabs}
\usepackage{longtable}
\usepackage{graphicx}
\usepackage{float}
\usepackage{xcolor}
\usepackage{tikz}
\usetikzlibrary{positioning,shapes.geometric,arrows.meta,calc}
\usepackage{enumitem}
\usepackage{ragged2e}
\usepackage[hidelinks]{hyperref}

% --- Encabezado y pie de página ---
\usepackage[myheadings]{fancyhdr}
\usepackage{lastpage}
\setlength{\headheight}{30pt}
\setlength{\footskip}{30pt}

% --- Interlineado ---
\linespread{1.2}
\sloppy

% =========================================================
% DATOS DEL DOCUMENTO
% =========================================================
\newcommand{\Institucion}{Instituto Tecnológico de Costa Rica}
\newcommand{\Campus}{Campus Tecnológico Central Cartago}
\newcommand{\DocumentoTitulo}{Entregable 09 -- Plan de Recursos}
\newcommand{\DocumentoSubtitulo}{}
\newcommand{\Curso}{Administración de proyectos}
\newcommand{\Profesor}{Alicia Marcela Salazar Hernández}
\newcommand{\FechaDocumento}{16 de junio de 2026}
\newcommand{\PeriodoAcademico}{I Semestre}
\newcommand{\AnioDocumento}{2026}

\newcommand{\AutoresPortada}{%
    Mauricio González Prendas: 2024143009\\
    Isaac Jiménez Blanco: 2024202804\\
    Tamara Robles Camacho: 2024099342%
}

% Datos que aparecen en el encabezado.
\newcommand{\DocHeaderLeft}{}
\newcommand{\DocHeaderRight}{Plan de Recursos}
\newcommand{\DocFooterRight}{Página~\thepage\ de~\pageref{LastPage}}

% Metadatos PDF.
\title{\DocumentoTitulo}
\author{\AutoresPortada}
\date{\FechaDocumento}

% =========================================================
% CONFIGURACIÓN DE TABLAS Y UTILIDADES
% =========================================================
\newcolumntype{Y}{>{\RaggedRight\arraybackslash}X}
\newcolumntype{L}[1]{>{\RaggedRight\arraybackslash}p{#1}}
\renewcommand{\arraystretch}{1.22}
\setlength{\LTpre}{0.5em}
\setlength{\LTpost}{0.5em}
\setcounter{tocdepth}{3}

% Imagen segura: si el archivo no existe, muestra un recuadro de marcador.
\newcommand{\SafeImage}[2][0.92\textwidth]{%
    \IfFileExists{#2}{%
        \includegraphics[width=#1]{#2}%
    }{%
        \fbox{%
            \begin{minipage}[c][4cm][c]{#1}
            \centering
            Imagen pendiente:\\
            \texttt{#2}
            \end{minipage}%
        }%
    }%
}

% Figura reutilizable.
\newcommand{\EvidenceFigure}[3]{%
    \begin{figure}[H]
    \centering
    \SafeImage{#1}
    \caption{#2}
    \label{#3}
    \end{figure}%
}

% =========================================================
% ENCABEZADO Y PIE
% =========================================================
\pagestyle{fancy}
\fancyhf{}
\fancyhead[L]{\DocHeaderLeft}
\fancyhead[R]{\DocHeaderRight}
\renewcommand{\headrulewidth}{0.4pt}
\renewcommand{\footrulewidth}{0.4pt}
\fancyfoot[R]{\DocFooterRight}

% Estilo equivalente para páginas horizontales.
\fancypagestyle{widefancy}{%
    \fancyhf{}%
    \fancyhead[L]{\DocHeaderLeft}%
    \fancyhead[R]{\DocHeaderRight}%
    \renewcommand{\headrulewidth}{0.4pt}%
    \renewcommand{\footrulewidth}{0.4pt}%
    \fancyfoot[R]{\DocFooterRight}%
}

% =========================================================
% PÁGINAS HORIZONTALES PARA TABLAS ANCHAS
% =========================================================
\makeatletter
\newcommand{\SyncPageBuilderDimensions}{%
    \global\vsize=\textheight
    \global\@colroom=\textheight
    \global\@colht=\textheight
}
\makeatother

\newenvironment{widepage}{%
    \clearpage
    \hypersetup{pageanchor=false}%
    \paperwidth=297mm
    \paperheight=210mm
    \pdfpagewidth=297mm
    \pdfpageheight=210mm
    \newgeometry{left=2.7cm,right=2.7cm,top=2.7cm,bottom=2.7cm}%
    \setlength{\textwidth}{243mm}%
    \setlength{\textheight}{156mm}%
    \setlength{\linewidth}{\textwidth}%
    \setlength{\columnwidth}{\textwidth}%
    \hsize=\textwidth
    \setlength{\headwidth}{\textwidth}%
    \SyncPageBuilderDimensions
    \pagestyle{widefancy}%
}{%
    \clearpage
    \paperwidth=210mm
    \paperheight=297mm
    \pdfpagewidth=210mm
    \pdfpageheight=297mm
    \restoregeometry
    \setlength{\textwidth}{156mm}%
    \setlength{\textheight}{243mm}%
    \setlength{\linewidth}{\textwidth}%
    \setlength{\columnwidth}{\textwidth}%
    \hsize=\textwidth
    \setlength{\headwidth}{\textwidth}%
    \SyncPageBuilderDimensions
    \hypersetup{pageanchor=true}%
    \pagestyle{fancy}%
}

% =========================================================
% PORTADA
% =========================================================
\newcommand{\MakeCover}{%
\begin{titlepage}
\thispagestyle{empty}
\centering

{\bfseries \huge \Institucion \par}
\vspace{0.25cm}
{\LARGE \Campus \par}

\vfill

{\huge \DocumentoTitulo \par}

\vfill

{\bfseries \LARGE Profesora: \par}
{\Large \Profesor \par}

\vfill

{\bfseries \LARGE Estudiantes: \par}
{\Large \AutoresPortada \par}

\vfill

{\bfseries\LARGE Fecha: \par}
{\Large \FechaDocumento \par}

\vfill

{\LARGE \PeriodoAcademico \par}
{\LARGE \AnioDocumento \par}

\end{titlepage}%
}

% =========================================================
% DOCUMENTO
% =========================================================
\begin{document}
\hypersetup{pageanchor=false}

% =========================
% Portada
% =========================
\MakeCover

% =========================
% Índice
% =========================
\pagenumbering{roman}
\pagestyle{empty}
\tableofcontents
\clearpage
\pagenumbering{arabic}
\hypersetup{pageanchor=true}
\pagestyle{fancy}

% =========================
% Contenido
% =========================
\begin{table}[H]
\centering
{\small
\begin{tabularx}{\textwidth}{L{4.5cm} Y}
\toprule
\textbf{Nombre del Proyecto} & Plataforma Universitaria Integrada \\
\midrule
\textbf{Organización} & Universidad Tecnológica La Mejor \\
\textbf{Tipo de Documento} & Plan de Recursos -- Entregable 9 \\
\textbf{Project Manager} & Tamara Robles Camacho \\
\textbf{Versión} & 1.0 \\
\textbf{Fecha} & Junio 2026 \\
\bottomrule
\end{tabularx}
}
\end{table}

\section{Recursos humanos}

El equipo del proyecto está conformado por cinco integrantes principales con dedicación completa durante la duración del proyecto, complementados por tres stakeholders clave con participación parcial en fases específicas.

\subsection{Equipo del proyecto}

En la siguiente tabla se muestran los roles de cada uno de los integrantes junto con su dedicación en el proyecto y las horas estimadas:

\begin{widepage}
{\small
\setlength{\tabcolsep}{3pt}
\begin{longtable}{L{4.3cm} L{4.6cm} L{2.2cm} L{3.1cm} L{5.1cm}}
\toprule
\textbf{Rol} & \textbf{Nombre / perfil} & \textbf{Dedicación} & \textbf{Horas estimadas} & \textbf{Fase(s) asignada(s)} \\
\midrule
\endfirsthead
\toprule
\textbf{Rol} & \textbf{Nombre / perfil} & \textbf{Dedicación} & \textbf{Horas estimadas} & \textbf{Fase(s) asignada(s)} \\
\midrule
\endhead
Project Manager & Tamara Robles Camacho & 100\% & 320 hrs & Todas las fases \\
Desarrollador Frontend \& QA 1 & Mauricio González Prendas & 100\% & 240 hrs & Ejecución / Pruebas \\
Desarrollador Frontend \& QA 2 & Carlos Sainz Arce & 100\% & 240h & Ejecución / Pruebas \\
Analista de Negocio \& Doc.1 & Isaac Jiménez Blanco & 100\% & 190 hrs & Inicio / Planificación / Cierre \\
Analista de Negocio \& Doc.2 & Jorge Abarca Quiros & 100\% & 190h & Inicio / Planificación / Cierre \\
Rector / Patrocinador & Dr. Roberto Mora Fallas & 5\% & 20 hrs & Inicio / Cierre \\
Directora Administrativa & Lic. Patricia Vargas Soto & 10\% & 30 hrs & Planificación / Ejecución \\
Director de TI & Ing. Carlos Brenes Mora & 10\% & 30 hrs & Planificación / Ejecución \\
\bottomrule
\end{longtable}
}
\end{widepage}

\subsection{Roles y responsabilidades clave}

Explicar los roles y sus responsabilidades es fundamental para llevar un orden en la elaboración del proyecto y evitar duplicación o mal información sobre responsabilidades, a continuación, se explican los roles y responsabilidades del equipo interno del proyecto:

\subsubsection{Project Manager (Tamara Robles Camacho)}

\begin{itemize}
    \item Responsable de la planificación, seguimiento y control del proyecto.
    \item Gestiona el registro de riesgos, cambios y comunicaciones con stakeholders.
    \item Elabora los entregables de gestión y coordina al equipo.
\end{itemize}

\subsubsection{Desarrollador Frontend \& QA 1 (Mauricio González Prendas)}

\begin{itemize}
    \item Implementa los módulos del portal estudiantil, docente y administrativo del sistema en React/HTML/CSS/JS.
    \item Realiza pruebas funcionales, de usabilidad y control de calidad del frontend.
    \item Responsable de la primera de la entrega del Entregable 18 (Frontend navegable).
\end{itemize}

\subsubsection{Desarrollador Frontend \& QA 2 (Carlos Sainz Arce)}

\begin{itemize}
    \item Implementa los módulos del portal financiero y ejecutivo del sistema en React/HTML/CSS/JS.
    \item Realiza, junto con el desarrollador 1, pruebas funcionales, de usabilidad y control de calidad del frontend.
    \item Responsable de la segunda parte de la entrega del Entregable 18 (Frontend navegable).
\end{itemize}

\subsubsection{Analista de negocio \& documentador 1 (Isaac Jiménez Blanco)}

\begin{itemize}
    \item Elabora la mitad de los entregables documentales de planificación y cierre.
    \item Gestiona la documentación de cronograma y comunicaciones.
    \item Apoya, junto con el analista de negocio 2, al PM en la coordinación de actividades de inicio y planificación.
\end{itemize}

\subsubsection{Analista de negocio \& documentador 2 (Jorge Abarca Quiros)}

\begin{itemize}
    \item Elabora la mitad de los entregables documentales de planificación y cierre.
    \item Gestiona la documentación de riesgos, costos y comunicaciones.
    \item Apoya, junto con él analista de negocio 1, al PM en la coordinación de actividades de inicio y planificación.
\end{itemize}

\subsection{Plan de capacitación}

Dado que el equipo es conformado por estudiantes universitarios con conocimiento técnico previo, se identifican las siguientes necesidades de nivelación:

\begin{itemize}
    \item MS Project: capacitación en línea (4 horas) para el responsable del cronograma.
    \item Metodología PMI/PMBOK: revisión guiada de verificación del PMI/PMBOK y plantillas (3 horas).
    \item React / Tailwind CSS: refuerzo técnico del desarrollador (autogestionado, 8 horas).
\end{itemize}

\section{Recursos tecnológicos}

A continuación, se detallan todas las herramientas de software, servicios en la nube y equipos de hardware requeridos para el desarrollo del proyecto.

\begin{widepage}
{\small
\setlength{\tabcolsep}{3pt}
\begin{longtable}{L{4cm} L{7cm} L{3cm} L{4cm} L{4.2cm}}
\toprule
\textbf{Recurso} & \textbf{Descripción / uso} & \textbf{Tipo} & \textbf{Costo} & \textbf{Responsable} \\
\midrule
\endfirsthead
\toprule
\textbf{Recurso} & \textbf{Descripción / uso} & \textbf{Tipo} & \textbf{Costo} & \textbf{Responsable} \\
\midrule
\endhead
VS Code & Editor de código principal para desarrollo frontend & Software & Gratuito & Mauricio \\
React + Vite & Framework y bundler para el frontend del sistema & Software & Gratuito & Mauricio \\
Tailwind CSS & Framework de estilos para la interfaz de usuario & Software & Gratuito & Mauricio \\
GitHub & Control de versiones y colaboración en código & Software & Gratuito & Mauricio \\
Figma & Diseño de wireframes y prototipos de interfaz & Software & Gratuito (plan free) & Mauricio \\
Microsoft Project & Elaboración del cronograma del proyecto & Software & Licencia comercial (\textasciitilde{}₡13,000/mes) & Isaac \\
Microsoft Word / Excel & Elaboración de entregables documentales & Software & Licencia Office 365 (\textasciitilde{}₡7,000/mes) & Todo el equipo \\
Cloudflare Workers & Hosting gratuito para despliegue del frontend demo & Servicio cloud & Gratuito & Mauricio \\
Correo y Teams & Comunicación formal del equipo en tiempo real, teams ofrece la posibilidad de reuniones virtuales del equipo y con stakeholders & Servicio & Gratuito & Las reuniones virtuales son responsables de Tamara y la comunicación de todo el equipo, \\
Computadoras personales (x5) & Equipos de trabajo personales de cada integrante & Hardware & Ya disponibles & Cada integrante \\
Conexión a Internet & Necesaria para reuniones virtuales y despliegue & Servicio & \textasciitilde{}₡25,000/mes c/u & Cada integrante \\
\bottomrule
\end{longtable}
}
\end{widepage}

\subsection{Justificación de herramientas seleccionadas}

En esta sección se justifican el porqué de la escogencia de dichas tecnologías y herramientas anteriormente mencionadas:

\begin{itemize}
    \item Se priorizan herramientas gratuitas o con licencia comercial simple para minimizar costos.
    \item React fue seleccionado por ser una de las tecnologías permitidas y por la experiencia previa del equipo.
    \item GitHub garantiza control de versiones y trazabilidad del código fuente.
    \item Cloudflare Workers permiten demostrar el frontend funcional sin costo de infraestructura.
\end{itemize}

\section{Recursos materiales}

Los recursos materiales necesarios para este proyecto son mínimos dado su carácter digital. Se identifican los siguientes materiales físicos de apoyo:

\begin{table}[H]
\centering
{\small
\setlength{\tabcolsep}{3pt}
\begin{tabularx}{\textwidth}{L{4cm} Y L{2.5cm} L{2.8cm}}
\toprule
\textbf{Material / recurso} & \textbf{Uso / descripción} & \textbf{Cantidad} & \textbf{Costo estimado} \\
\midrule
Impresiones de documentos & Entregables formales impresos para revisión presencial & \textgreater{}80 páginas & ₡9,000 \\
Cuadernos / libretas & Apuntes de reuniones y planificación manual & 5 unidades & ₡9,500 \\
Bolígrafos y marcadores & Uso en sesiones presenciales y pizarras & Set básico & ₡5,000 \\
Pizarra (uso en las oficinas) & Sesiones de planificación y diseño colaborativo & 1 & ₡0 \\
USB / almacenamiento externo & Respaldo de archivos del proyecto & 1 por integrante & ₡10,000 \\
Carpetas y archivadores & Organización física de entregables impresos & 5 unidades & ₡6,500 \\
\midrule
\textbf{TOTAL RECURSOS MATERIALES} & \textbf{--} & \multicolumn{2}{l}{\textbf{₡40,000}} \\
\bottomrule
\end{tabularx}
}
\end{table}

\section{Resumen de recursos del proyecto}

La siguiente tabla integra todos los recursos identificados por categoría:

\begin{table}[H]
\centering
{\small
\setlength{\tabcolsep}{3pt}
\begin{tabularx}{\textwidth}{L{5.5cm} Y L{4.3cm}}
\toprule
\textbf{Categoría} & \textbf{Descripción} & \textbf{Costo estimado} \\
\midrule
Recursos Humanos & 5 integrantes + 3 stakeholders parciales & ₡14,000,000 (mano de obra total) \\
Recursos Tecnológicos -- Software & 10 herramientas identificadas, 3 con pago & ₡45,000 (licencias pagas) \\
Recursos Tecnológicos -- Hardware & 5 computadoras personales & Ya disponibles (sin costo adicional) \\
Recursos Tecnológicos -- Servicios & Internet, hosting, reuniones virtuales & ₡75,000 (3 meses x 5 personas) \\
Recursos Materiales & Impresiones, papelería, USBs & ₡30,000 \\
\midrule
\textbf{INVERSIÓN TOTAL ESTIMADA DEL PROYECTO} & \multicolumn{2}{l}{\textbf{₡16,100,000}} \\
\bottomrule
\end{tabularx}
}
\end{table}

\end{document}
